\section{Projects}
\label{sec:cses-projects}

\problemstatement{There are $n$ projects, each with a start day, an end
day, and a reward. Attending a project occupies all days from its start to
its end inclusive, so two attended projects may not overlap. Maximise the
total reward.}

\begin{patternbox}
\emph{``Non-overlapping intervals, maximise total value''} is
\textbf{weighted interval scheduling}. Sort projects by end day; for each,
either skip it or take it and jump to the best solution that finishes
before its start -- found by binary search.
\end{patternbox}

\subsection*{State and transition}

Sort the projects by end day. Let $\textit{dp}[i]$ be the maximum reward
using only the first $i$ projects (in sorted order). For project $i$ with
start $s_i$ and reward $p_i$, either skip it ($\textit{dp}[i-1]$) or take
it: add $p_i$ to the best reward achievable among projects that end
\emph{strictly before} $s_i$. Let $j$ be the number of projects whose end
day $<s_i$ (found by binary search on the sorted end days):
\[
  \textit{dp}[i]=\max\big(\textit{dp}[i-1],\ p_i+\textit{dp}[j]\big).
\]

\begin{ideabox}[Sort by End, Binary-Search the Predecessor]
Sorting by end day means that when we consider a project, every compatible
earlier project has already been folded into some $\textit{dp}[j]$. Binary
search finds the latest non-conflicting project in $O(\log n)$, turning an
$O(n^2)$ scan into $O(n\log n)$.
\end{ideabox}

\subsection*{Algorithm}

\begin{enumerate}
  \item Sort projects by end day; let $\textit{ends}[]$ be their end
        days.
  \item $\textit{dp}[0]=0$. For $i=1\ldots n$: binary-search $j$ = number
        of projects with end $<s_i$; set
        $\textit{dp}[i]=\max(\textit{dp}[i-1], p_i+\textit{dp}[j])$.
  \item Answer $\textit{dp}[n]$.
\end{enumerate}

\subsection*{Dry run}

\dryrun{Projects (start, end, reward): $(2,4,7),(1,2,4),(5,6,3)$.}

\begin{itemize}
  \item Sorted by end: $(1,2,4),(2,4,7),(5,6,3)$.
  \item $\textit{dp}[1]=4$. Project $(2,4,7)$: projects ending $<2$ is
        the first (end $2$? no, need $<2$, so $j=0$):
        $\textit{dp}[2]=\max(4,7+0)=7$.
  \item Project $(5,6,3)$: projects ending $<5$ are both, $j=2$:
        $\textit{dp}[3]=\max(7,3+7)=10$. Answer $10$.
\end{itemize}

\subsection*{C++ implementation}

\begin{lstlisting}
#include <bits/stdc++.h>
using namespace std;

int main() {
    int n;
    cin >> n;
    vector<array<long long,3>> p(n);           // {end, start, reward}
    for (int i = 0; i < n; i++) {
        long long a, b, x;
        cin >> a >> b >> x;
        p[i] = {b, a, x};
    }
    sort(p.begin(), p.end());                  // by end day

    vector<long long> ends(n);
    for (int i = 0; i < n; i++) ends[i] = p[i][0];

    vector<long long> dp(n + 1, 0);
    for (int i = 1; i <= n; i++) {
        long long start = p[i-1][1], reward = p[i-1][2];
        // number of projects with end < start
        int j = lower_bound(ends.begin(), ends.end(), start) - ends.begin();
        dp[i] = max(dp[i-1], reward + dp[j]);
    }
    cout << dp[n] << "\n";
    return 0;
}
\end{lstlisting}

\begin{complexitybox}
\textbf{Time Complexity.} $O(n\log n)$ -- sorting plus a binary search per
project. \textbf{Space Complexity.} $O(n)$.
\end{complexitybox}

\begin{takeawaybox}
Weighted interval scheduling = sort by end, then ``skip or take with a
binary-searched predecessor.'' It generalises the unweighted greedy
(which just maximises count) to rewards, where greed fails and DP is
required.
\end{takeawaybox}
